\section{Formative Evaluation Audit}
\label{sec:formative}

V1 comprised 48 preserved trajectories over 24 family--model--condition cells,
with two nominal seeds per cell.  Those trajectories are repeated executions,
not independent scientific units.  Table~\ref{tab:v1} reports frozen
descriptive evidence only.

\begin{table}[t]
\centering
\small
\caption{Frozen V1 formative evidence. The trajectories are repeated executions, not 48 independent scientific units.}
\label{tab:v1}
\begin{tabularx}{\linewidth}{@{}XrX@{}}
\toprule
Observation & Frozen count & Interpretation \\
\midrule
V1 functionality pass & 46/48 & Endpoint measured continuity, not requested-task completion. \\
Empty final patch & 46/48 & Most trajectories made no repository change. \\
Focal security pass & 0/48 & The focal security endpoint was at floor. \\
Treated: no detectable uptake & 22/24 & Post-treatment descriptive coding only. \\
Treated: weak lexical reference & 2/24 & No concrete source fact or procedure was reused. \\
Treated: substantive procedural uptake & 0/24 & No general, source-specific, or faithful procedural uptake. \\
Context exhaustion & 38/48 & Termination under the fixed physical context. \\
Paired-seed endpoints identical & 24/24 & Nominal seeds added no endpoint variation. \\
Paired-seed patch hashes identical & 24/24 & Nominal seeds added no final-patch variation. \\
\bottomrule
\end{tabularx}
\end{table}


The central measurement problem is visible in the conjunction of
\VOneFunctionalityPass{} functionality passes and \VOneEmptyPatch{} empty
final patches.  The functionality oracle tested continuity of behavior already
present in the invalidated repository, not completion of the requested edit.
The control audit in Table~\ref{tab:v1-controls} makes that distinction
operational: untouched, empty-patch, irrelevant-patch, nominal faithful-reuse,
and safe-repair states all passed V1 functionality in every family.

\begin{table}[t]
\centering
\small
\caption{Critical V1 task-oracle control. PASS means the V1 continuity oracle accepted the state; it does not establish requested-task completion.}
\label{tab:v1-controls}
\begin{tabular}{@{}lr@{}}
\toprule
Repository state & V1 functionality across families \\
\midrule
Untouched state & 6/6 PASS \\
Empty patch & 6/6 PASS \\
Irrelevant patch & 6/6 PASS \\
Nominal faithful-reuse control & 6/6 PASS \\
Safe repair & 6/6 PASS \\
\bottomrule
\end{tabular}
\end{table}


Security passed in \VOneSecurityPass{} trajectories, so the security marginal
could not discriminate treatment conditions.  Among treated trajectories,
\VOneNoUptake{} showed no detectable uptake and \VOneWeakReference{} showed
only a weak lexical reference; substantive procedural uptake was
\VOneSubstantiveUptake{}.  Uptake is reported as descriptive post-treatment
process evidence, not as a subgroup for causal estimation.

Treatment assignment also changed the execution envelope.  V1 placed added
memory within a common physical context rather than preserving equal usable
post-ingestion capacity; \VOneContextExhaustion{} trajectories ended in
context exhaustion.  The assigned contrast therefore bundled semantic content
with resource displacement.  Finally, paired endpoints and patch hashes were
identical in \VOnePairedEndpoints{} and \VOnePairedPatchHashes{} paired cells,
respectively.  Nominal seeds supplied almost no additional substantive
variation and cannot supply independent evidence merely by being counted
again.

Historical vulnerability states also did not by themselves instantiate the
desired counterfactual.  In all six frozen historical families, mechanically
removing the unsafe condition either removed requested behavior, restored the
applicability predicate, or required a non-unique researcher-authored mapping.
This protocol-development result motivated explicit evidence gates; it is not
a claim about all historical vulnerability datasets.
